\documentclass[aspectratio=43]{beamer}
\usepackage[english]{babel}
\usepackage{amsthm}
\usepackage{mathtools}
% \usepackage{physics}
\usepackage{calligra}
\usepackage{csquotes}
%\usepackage{tensor}
\usepackage[thicklines]{cancel}
\usepackage{tcolorbox}
\usepackage{pstricks}
%\usepackage[backend=biber, bibstyle=nature, sorting=nty, citestyle=numeric-comp]{biblatex} %Custom bibliography
%    \addbibresource{bib.bib} %Load references


\DeclareMathAlphabet{\mathcalligra}{T1}{calligra}{m}{n}
\DeclareFontShape{T1}{calligra}{m}{n}{<->s*[2.2]callig15}{}
\newcommand{\scriptr}{\mathcalligra{r}\,}
\newcommand{\boldscriptr}{\pmb{\mathcalligra{r}}\,}
\def\rc{\scriptr}
\def\brc{\boldscriptr}
\def\hrc{\hat\brc}
\newcommand{\ie}{\emph{i.e.}} %id est
\newcommand{\eg}{\emph{e.g.}} %exempli gratia
\newcommand{\rtd}[1]{\ensuremath{\left\lfloor #1 \right\rfloor}}
\newcommand{\dirac}[1]{\ensuremath{\delta \left( #1 \right)}}
\newcommand{\diract}[1]{\ensuremath{\delta^3 \left( #1 \right)}}
\newcommand{\e}{\ensuremath{\epsilon_0}}
\newcommand{\m}{\ensuremath{\mu_0}}
\newcommand{\V}{\ensuremath{\mathcal{V}}}
\newcommand{\prnt}[1]{\ensuremath{\left(#1\right)}} %parentheses
\newcommand{\colch}[1]{\ensuremath{\left[#1\right]}} %square brackets
\newcommand{\chave}[1]{\ensuremath{\left\{#1\right\}}}  %curly brackets

\useoutertheme{infolines}
\useinnertheme{rectangles}
\usefonttheme{professionalfonts}


\definecolor{orange}{HTML}{f28165}
\definecolor{gray}{HTML}{303030}
\definecolor{yellow}{HTML}{f0be52}
\definecolor{lightorange}{HTML}{f19e58}

\renewcommand{\CancelColor}{\color{orange}}

\makeatletter
\newcommand{\mybox}[1]{%
  \setbox0=\hbox{#1}%
  \setlength{\@tempdima}{\dimexpr\wd0+13pt}%
  \begin{tcolorbox}[colback=orange,colframe=orange,boxrule=0.5pt,arc=4pt,
      left=6pt,right=6pt,top=6pt,bottom=6pt,boxsep=0pt,width=\@tempdima]
    \textcolor{white}{#1}
  \end{tcolorbox}
}
\makeatother

\usecolortheme[named=orange]{structure}
\usecolortheme{sidebartab}
\usecolortheme{orchid}
\usecolortheme{whale}
\setbeamercolor{alerted text}{fg=yellow}
\setbeamercolor{block title alerted}{bg=alerted text.fg!90!black}
\setbeamercolor{block title example}{bg=lightorange!60!black}
\setbeamercolor{background canvas}{bg=gray}
\setbeamercolor{normal text}{bg=gray,fg=white}

\setbeamertemplate{footline}
        {
      \leavevmode%
      \hbox{%
      \begin{beamercolorbox}[wd=.333333\paperwidth,ht=2.25ex,dp=1ex,center]{author in head/foot}%
        \usebeamerfont{author in head/foot}\insertshortauthor~~(\insertshortinstitute)
      \end{beamercolorbox}%
      \begin{beamercolorbox}[wd=.333333\paperwidth,ht=2.25ex,dp=1ex,center]{title in head/foot}%
        \usebeamerfont{title in head/foot}\insertshorttitle
      \end{beamercolorbox}%
      \begin{beamercolorbox}[wd=.333333\paperwidth,ht=2.25ex,dp=1ex,center]{date in head/foot}%
        \usebeamerfont{date in head/foot}\hfill\insertframenumber/\inserttotalframenumber%\hspace*{2em}

    %#turning the next line into a comment, erases the frame numbers
        %\insertframenumber{} / \inserttotalframenumber\hspace*{2ex} 

      \end{beamercolorbox}}%
      \vskip0pt%
    }


\setbeamertemplate{blocks}[rectangle]
\setbeamercovered{dynamic}

\setbeamertemplate{section page}
{
	\begin{centering}
		\begin{beamercolorbox}[sep=27pt,center]{part title}
			\usebeamerfont{section title}\insertsection\par
			\usebeamerfont{subsection title}\insertsubsection\par
		\end{beamercolorbox}
	\end{centering}
}

%\setbeamertemplate{subsection page}
%{
%	\begin{centering}
%		\begin{beamercolorbox}[sep=12pt,center]{part title}
%			\usebeamerfont{subsection title}\insertsubsection\par
%		\end{beamercolorbox}
%	\end{centering}
%}

\newcommand{\hlight}[1]{\colorbox{violet!50}{#1}}
\newcommand{\hlighta}[1]{\colorbox{red!50}{#1}}
\title{Fun with Ubuntu} %->->->->-> Check hyperref title <-<-<-<-<-
\subtitle{Just Dive into It!}
\author[S. T. Miah]{Shajeda (Mahee) T. Miah}
\institute[HGES]{
    Fifth Grade%
    \\
    Hickory Grove Elementary School%
    \\
    ~\\
    2514 W Hickory Grove Rd, Dunlap, IL 61525\\
    UNITED STATES
} %You can change the Institution if you are from somewhere else
\date{\today}
%\logo{\includegraphics[width= 0.2\textwidth]{images/a-logo.png}}

\begin{document}
    
    \frame{\titlepage}
    
    \begin{frame}{Outline}
        \tableofcontents
    \end{frame}

	\section{Ubuntu Overview}
			
	\begin{frame}{Ubuntu Overview}

\begin{block}{What is Ubuntu}
	\begin{itemize}
	\item  Ubuntu is a Linux operating system. 
	\item This operating system is for computers, smartphones plus network servers.
	\item Ubuntu is developed by a company called Canonical Ltd, which was based off of a UK company.
	\item Ubuntu is released (officially) in three editions (editions below.)
	\\
	~\\
	\item Desktop
	\item Server
	\item Core
	\end{itemize}		
\end{block}
\end{frame}		
	
\begin{frame}{Features of Ubuntu}
\begin{itemize}
\item Desktop version supports of normal software 

\item Supports office like ``LibreOffice"

\item Ubuntu has a built-in email software called ``Thunderbird," you can access emails that are in Gmail, Hotmail, and Exchange

\item There is a host for free applications for users of Ubuntu to view and edit.

\item In Ubuntu, there are also applications to manage vidoes and lets the user share vidoes as well.

\item It is easy to find satisfaction on ubuntu with the ``smart searching facitlity"

\item The best feature is that Ubuntu is a free operating system, plus it is backed by a ``huge open source community."


\end{itemize}


\end{frame}	

\section{What is the Actual Purpose of Ubuntu}
\begin{frame}{What is the Actual Purpose of Ubuntu?}
	\begin{itemize}
		\item The actual purpose of using Ubuntu is down below after this line.
		\item Now the main thing now is what is the exact purpose of Ubuntu, the main purpose of Ubuntu is that it is to give the user seemless use of the desktop like you would use for windows.
	\item You have better control of the operating system and the computer that you are using than when you are working with windows.
	\end{itemize}
\end{frame}



	\section{What "Flavours" does Ubuntu have }
	\begin{frame}{What "Flavours" does Ubuntu have?}
	\begin{block}{}
	\begin{itemize}
	\item Ubuntu comes in many different versions and not the type of 20.04 version. I meant like some other versions like Kubuntu or Xubuntu. The list of the most popular ones are down below.
	\item Ubuntu GNOME (This is the regular Ubuntu. Ubuntu is the main plus the most popular flavour of Ubuntu. This Ubuntu runs on GNOME desktop environment. )
	\item Kubuntu (Kubuntu is known as the official flavour of the Ubuntu Operating system. Kubuntu runs on KDE instead of GNOME.)
	\end{itemize}
	\end{block}
	\end{frame}
	
	
\begin{frame}{Continuation of other page.}
\begin{block}{}
\begin{itemize}

\item Ubuntu Desktop (This flavour of Ubuntu is not exactly a flavour but it is one Ubuntu and it is also one of the three editions of Ubuntu (as said in slide two.) So I would consider it still an Ubuntu.)
\item Xubuntu (Xubuntu is one flavour of Ubuntu, Xubuntu works quite well with old computer hardware so it works well with an old computer, and it's installed desktop environment is XFCE.)
\item Lubuntu (Lubuntu is a lightweight Linux distribution that is based of of Ubuntu. This flavour of Ubuntu runs on the LXqT desktop environment and as most of the ones that I mentioned here Lubuntu does not run on GNOME desktop environment.)

\end{itemize}
\end{block}
\end{frame}


\section{What are Some Commands on Ubuntu Terminal}
\begin{frame}{What are Some Commands on Ubuntu Terminal?}
\begin{itemize}

\item So if I mentioned about terminal, terminal is a command line where you would want to open any files or go to any app like firefox web browser or texmaker (which is where I am making my presentation.)
\item There are some commands you would need to know to do these type of things.
\item One command you could do is push something to github.
\item Github is a cloud-based service for developers to manage and store their code, you can also make repositories in github.
\item Then github comes into this because this slide will teach you how to push stuff from one README file that you might have in your homefile to your github page.
\item The commands will be on the next page.

\end{itemize}
\end{frame}

\begin{frame}{The Commands for Pushing the Edits you Make to Your Github Page README File From Terminal.}
\begin{itemize}
\item The first thing you need to type is git status to see what you modified lately.
\item Then the second thing you need to do is git add . for if you want to add anything.
\item Next the third step is to type git commit -m then "some meaning full message" to leave in your github page to show that you modified the file.
\item the final step is to type git push to push those edits to your github page, after that you just type in your github username and github password.
\item And that is it.
\end{itemize}
\end{frame}



\begin{frame}{Another few Commands That you Could use on Terminal}
\begin{itemize}

\item There are some other commands that you could use in terminal.
\item These commands help you see that if you have any programming language in your Ubuntu like python.
\item The commands are down below this line.
\\
~\\ 
\item Command number 1: python -v (has to be lowercase.)
\item Command number 2: python -V (has to be capital.)
\item Command number 3: python3 -V (You could see if there is Python3 on your computer terminal as well.)
\end{itemize}
\end{frame}

\section{What are Some Applications (apps)that you have on Ubuntu}
\begin{frame}{What are Some Applications (apps)that you have on Ubuntu?}
\begin{itemize}
\item There are lots of apps that you would have on Ubuntu more than you might have on windows.
\item Like for windows unless you download terminal you would have it staight away once you download Ubuntu.
\item Also you would have Mahjohhng which is a almost like a puzzle that you have to solve so it helps with critical thinking.
\item Then for windows you would have mainly google as the seach engine, but in Ubuntu the web browser is Firefox.
\item Then last but not least that you would not have in windows unless you download this would be Rythmbox which you can save music that you save from the web and you can listen to it any time you want if you are working.
\item There are more apps but I dicided to go with these few apps because when I downloaded Ubuntu 20.04 these apps were already there.
\end{itemize}
\end{frame}

\section{Virtual Machines on Ubuntu}
\begin{frame}{Virtual Machines on Ubuntu}
\begin{itemize}
\item Ubuntu can also be installed as virtual machines.
\item Some of these vitual machines are like Oracle VirtualBox.
\item There are also Virtual machines. 
\item The other ones are down below.
\\
~\\
\item Microsoft hyper-V
\item VMWare Workstation
\\
~\\
\item So Virtual Machines are safe and the safest one is one that I did not mention there so I am mentioning now.
\item The safest virtual machine is virtual machines 101.
\item So let's say that your virtual machine gets hacked your "host" which is the operating system you are running which is Ubuntu is still safe.
\item Make sure that you still have an operating system running that way if your VM gets hacked then the operating system is still safe.
\end{itemize}
\end{frame}
\section{How Many Versions are There of Ubuntu}
\begin{frame}{How Many Versions are There of Ubuntu?}
\begin{itemize}
\item There are many versions of Ubuntu, there is Ubuntu 20.04 or Ubuntu 16.04. Mainly, there are many versions of ubuntu.
\item The most popular versions of Ubuntu are the ones down below.
\\
~\\
\item Ubuntu 20.04 (This is the operating system that my macbook pro has.)
\item Ubuntu 18.04 (Getting followed by Ubuntu 20.04)
\item Ubuntu 16.04 (Released in 2016)
\item Ubuntu 14.04 ( Released in 2014)
\item These are some of the most popular versions so for now I will call it that there are four versions for now, but there are many more than just these. Those are just the most popular versions that I mentioned. 
\end{itemize}
\end{frame}

	
\section{The End}
    \begin{frame}{ }
        \centering
            \Huge\bfseries
        \textcolor{orange}{The End -- Thank you!}
    \end{frame}


\end{document}
